% Core cluster: corrected grids, filtering, growth-based state bounds,
% unknown-growth schedule, bit complexity, and lower bounds.
% Self-contained fragment; no preamble.  Uses \R \Z \E \norm \poly \OPT.

\subsection{Setting and the conditioning parameter}\label{core:sec:setting}

Let $X=\prod_{i=1}^nX_i$, where $X_i=[l_i,u_i]$ for $i\in I_C$ (continuous
coordinates) and $X_i=[l_i,u_i]\cap\Z$ for $i\in I_Z$ (integer coordinates).
Integer endpoints are rounded inward, empty domains are rejected, and fixed
coordinates are substituted, so that $n\ge1$ and $l_i<u_i$ for all $i$, with
$l_i,u_i\in\Z$ for $i\in I_Z$. Write $s_i=u_i-l_i$, $s=\max_is_i$, and
$\bar X=\prod_i[l_i,u_i]$ for the continuous hull. The objective is
\[
F(x)=\sum_{a\in\mathcal A}f_a(x_{S_a}),\qquad \OPT=\min_{x\in X}F(x),
\]
with each $f_a$ defined on the continuous hull of its scope. A tree
decomposition $(T,(B_t)_{t\in T})$ is supplied: every scope $S_a$ lies in
some bag, and for every $i$ the bags containing $i$ form a nonempty subtree.
Let $N=|T|$ and $p=\max_t|B_t|$. Each factor is assigned to one bag containing
its scope, and each coordinate $i$ to one \emph{home bag} containing it.

\begin{definition}[Coordinate curvature]\label{core:def:curv}
Numbers $L_1,\dots,L_n\ge0$ are \emph{coordinate curvature bounds} if, for
every $i$ and every $x\in\bar X$, the function
$t\mapsto F(x_1,\dots,x_{i-1},t,x_{i+1},\dots,x_n)-L_it^2/2$ is concave on
$[l_i,u_i]$. Put $P=\{i:L_i>0\}$, $n_P=|P|$, $L=\max_iL_i$, and
\[
\norm{d}_{\mathbf L}^2=\sum_{i\in P}L_id_i^2\qquad(d\in\R^n).
\]
\end{definition}

The condition concerns the sum $F$ and the whole continuous hull, including
fractional values of integer coordinates. For a quadratic with Hessian $H$ the
smallest valid bounds are $L_i=\max\{0,H_{ii}\}$. Larger numbers are also
valid; in particular the common bound $L_i=L$ is valid.

\begin{definition}[Growth and conditioning]\label{core:def:growth}
The instance has \emph{weighted growth} with constant $\gamma>0$ if there is
a global minimizer $x^*$ with
\begin{equation}\label{core:eq:growth}
F(x)-\OPT\ge\gamma\norm{x-x^*}_{\mathbf L}^2\qquad(x\in X).
\end{equation}
Let $\gamma^*$ be the largest such constant (the set of valid constants is
closed) and put $\bar\kappa=\max\{1,1/\gamma^*\}$; put $\bar\kappa=\infty$ if
no $\gamma>0$ exists. The usual Euclidean growth is
$F(x)-\OPT\ge g\norm{x-x^*}^2$, with $\kappa=\max\{1,L/g\}$.
\end{definition}

\begin{remark}\label{core:rem:kappa}
(i) Since $\norm{d}_{\mathbf L}^2\le L\norm d^2$, Euclidean growth with
constant $g$ implies weighted growth with $\gamma\ge g/L$; hence
$\bar\kappa\le\kappa$. All results below that are stated with $\bar\kappa$
therefore hold with $\kappa$.
(ii) Replacing a continuous coordinate $x_i$ by $x_i=\alpha y_i$ multiplies
$L_i$ by $\alpha^2$ and leaves $L_i(x_i-x_i^*)^2$ unchanged; $\bar\kappa$ is
invariant under such rescalings, whereas $\kappa$ is not. For a separable
quadratic $\sum_i\frac{L_i}2(x_i-c_i)^2$ with interior minimizer,
$\bar\kappa=2$ while $\kappa=2\max_iL_i/\min_iL_i$.
(iii) Coordinates with $L_i=0$ are concave in their own direction; they need
neither growth nor localization. Weighted growth makes $x^*_P$ unique but
allows global minimizers that differ outside $P$. For example, binary
variables whose squares are written with nonpositive coefficients are free.
(iv) Growth (and hence uniqueness of $x^*_P$) is used only to bound running
time. No validity statement below uses it.
\end{remark}

\subsection{Corrected grids give continuous lower bounds}\label{core:sec:grid}

A \emph{subbox} is $B=\prod_iB_i\subseteq X$ with $B_i=[l_i',u_i']$ or
$[l_i',u_i']\cap\Z$ of the same type as $X_i$, with $l_i'\le u_i'$ and
integer endpoints for $i\in I_Z$. A \emph{grid} for $B$ is $G=\prod_iG_i$ with
finite $G_i\subseteq B_i$ containing $l_i'$ and $u_i'$. Consecutive nodes of
$G_i$ form the \emph{grid intervals} of coordinate $i$.

\begin{definition}[Corrections]\label{core:def:corr}
For $v\in G_i$ let $w_i(v)$ be the largest length of a grid interval with
endpoint $v$, where for $i\in I_Z$ intervals of length one are ignored;
$w_i(v)=0$ if no interval remains. Put
\[
d_i(v)=\tfrac{L_i}{8}w_i(v)^2,\quad D(y)=\sum_id_i(y_i),\quad Q=F-D,\quad
\beta=\min_{y\in G}Q(y),\quad
m_i(v)=\min\{Q(y):y\in G,\ y_i=v\}.
\]
\end{definition}

\begin{lemma}[Interpolation]\label{core:lem:interp}
For every $x\in B$ there is a random $Y\in G$ with independent coordinates,
$\E Y=x$, and $\E Q(Y)\le F(x)$, such that $Y_i=x_i$ if $x_i\in G_i$, and
otherwise $Y_i$ takes the two endpoints of the grid interval containing
$x_i$. Consequently:
\begin{enumerate}
\item[(a)] $\beta\le\min_BF$;
\item[(b)] if $a<a'$ are consecutive in $G_i$, then
$\min\{m_i(a),m_i(a')\}\le F(x)$ for every $x\in B$ with $x_i\in[a,a']$;
\item[(c)] if $y$ minimizes $Q$ over $G$, then $F(y)-\beta=D(y)\ge0$.
\end{enumerate}
\end{lemma}

\begin{proof}
If $x_i\notin G_i$, let $a<x_i<a'$ be consecutive nodes and let $Y_i=a'$ with
probability $(x_i-a)/(a'-a)$ and $Y_i=a$ otherwise; then $\E Y_i=x_i$ and
$\operatorname{Var}Y_i=(x_i-a)(a'-x_i)\le(a'-a)^2/4$. If $x_i\in G_i$, let
$Y_i=x_i$. Take the $Y_i$ independent and put
$Y^{(k)}=(Y_1,\dots,Y_k,x_{k+1},\dots,x_n)\in\bar X$.
Fix $Y_1,\dots,Y_{k-1}$. The function
$\varphi(t)=F(Y_1,\dots,Y_{k-1},t,x_{k+1},\dots,x_n)-L_kt^2/2$ is concave on
$[l_k,u_k]$, so Jensen's inequality gives $\E\varphi(Y_k)\le\varphi(x_k)$,
that is,
\[
\E\bigl[F(Y^{(k)})\,\big|\,Y_1,\dots,Y_{k-1}\bigr]
\le F(Y^{(k-1)})+\tfrac{L_k}2\operatorname{Var}Y_k .
\]
Taking expectations and summing over $k$ gives
$\E F(Y)\le F(x)+\sum_k\frac{L_k}2\operatorname{Var}Y_k$.

If $\operatorname{Var}Y_k>0$ then $a<x_k<a'$. For $k\in I_Z$ the numbers
$x_k,a,a'$ are integers, so $a'-a\ge2$ and the interval is not ignored. In all
cases $w_k(a),w_k(a')\ge a'-a$, hence
$\E d_k(Y_k)\ge L_k(a'-a)^2/8\ge\frac{L_k}2\operatorname{Var}Y_k$. If
$\operatorname{Var}Y_k=0$ then $\E d_k(Y_k)\ge0$. Therefore
$\E D(Y)\ge\sum_k\frac{L_k}2\operatorname{Var}Y_k$ and $\E Q(Y)\le F(x)$.

(a) $\beta\le\E Q(Y)\le F(x)$. (b) If $x_i\in[a,a']$ then $Y_i\in\{a,a'\}$
surely, so $Q(Y)\ge m_i(Y_i)\ge\min\{m_i(a),m_i(a')\}$ and the claim follows by
taking expectations. (c) holds by definition and $D\ge0$.
\end{proof}

For $i\notin P$ we always use $G_i=\{l_i',u_i'\}$; then $d_i\equiv0$.

\subsection{One tree computation gives all coordinate bounds}\label{core:sec:dp}

Let $q_t$ be the sum of the factors assigned to bag $t$ minus the corrections
$d_i$ of the coordinates whose home bag is $t$; thus $Q=\sum_tq_t$ on $G$. For
adjacent $t,u$ put $S_{tu}=B_t\cap B_u$ and define messages
\[
M_{t\to u}(x_{S_{tu}})=\min_{x_{B_t\setminus S_{tu}}}
\Bigl[q_t(x_{B_t})+\sum_{v\in\mathcal N(t)\setminus\{u\}}M_{v\to t}(x_{S_{vt}})\Bigr],
\]
all variables ranging over their grids (an empty separator has one state).
Let $A_t=q_t+\sum_{v\in\mathcal N(t)}M_{v\to t}$.

\begin{proposition}[Tree dynamic programming]\label{core:prop:dp}
For every bag $t$, $A_t(x_{B_t})=\min\{Q(y):y\in G,\ y_{B_t}=x_{B_t}\}$. Hence
$\min A_t=\beta$, and $m_i(v)$ is the minimum of $A_t$ over $x_i=v$ for any bag
$t\ni i$. With $K=\max_i|G_i|$, all messages, $\beta$, a minimizer, and all
$m_i$ are obtained with $O\bigl((|\mathcal A|+n+N)K^p\bigr)$ additions and
comparisons on table entries, plus the evaluation of each factor at its at
most $K^{|S_a|}$ grid points, plus $O(p)$ index work per table entry.
\end{proposition}

\begin{proof}
For an oriented edge $t\to u$, let $T_{t\to u}$ be the component of
$T\setminus\{tu\}$ containing $t$, and let $V_{t\to u}$ be the set of
coordinates occurring in bags of $T_{t\to u}$ but not in $S_{tu}$. We show by
induction on $|T_{t\to u}|$ that $M_{t\to u}(x_{S_{tu}})$ is the minimum of
$\sum_{r\in T_{t\to u}}q_r$ over $x_{V_{t\to u}}$. The trees $T_{v\to t}$,
$v\in\mathcal N(t)\setminus\{u\}$, partition $T_{t\to u}\setminus\{t\}$. If a
coordinate occurs in a bag of $T_{v\to t}$ and in a bag outside it, running
intersection puts it in $B_v$ and $B_t$, i.e., in $S_{vt}$. Hence the sets
$V_{v\to t}$ are pairwise disjoint, disjoint from $B_t$, and
$V_{t\to u}=(B_t\setminus S_{tu})\cup\bigcup_vV_{v\to t}$ is a disjoint union.
The sum $\sum_{r\in T_{v\to t}}q_r$ depends only on $x_{V_{v\to t}}$ and
$x_{S_{vt}}$, so the minimization splits, and by induction each inner minimum
is $M_{v\to t}$. This proves the claim; the same argument with all neighbours
of $t$ gives the formula for $A_t$, since $B_t$ and the sets $V_{v\to t}$,
$v\in\mathcal N(t)$, partition all coordinates.

Compute messages towards a root and then away from it. At bag $t$ form the
table $q_t$ (one addition per assigned factor or correction per entry), add
all incoming messages once to obtain $A_t$, and obtain each outgoing message
from $A_t-M_{u\to t}$ by one pass with a comparison per entry; exact
arithmetic makes the subtraction legitimate. Each bag table has at most $K^p$
entries and $\sum_t|\mathcal N(t)|=2(N-1)$, which gives the bound. A minimizer
is recovered by choosing, from the root outwards, a minimizing entry of $A_t$
consistent with the labels already fixed; consistency of all $A_t$ at the value
$\beta$ follows from the first claim.
\end{proof}

This is the classical nonserial dynamic programming or junction-tree
computation \cite{BerteleBrioschi1972,Dechter1999}; the only new element is
the unary correction that turns it into a lower bound for the continuous
problem.

\subsection{Filtering and path certificates}\label{core:sec:filter}

\begin{definition}[Filtering step]\label{core:def:filter}
Let $G$ be a grid for $B$ and $U\in\R$. For $i\in P$, \emph{retain} a grid
interval $[a,a']$ of coordinate $i$ if $\min\{m_i(a),m_i(a')\}\le U$, and let
$B_i'$ be the hull of the retained intervals (intersected with $\Z$ for
$i\in I_Z$). For $i\notin P$ let $B_i'=B_i$.
\end{definition}

\begin{proposition}[Safe filtering]\label{core:prop:filter}
Every $x\in B\setminus B'$ satisfies $F(x)>U$. Consequently, if
$U\ge\min_BF$, then $B'$ contains every $x\in B$ with $F(x)\le U$, in
particular every minimizer of $F$ over $B$. If moreover $y$ minimizes $Q$ and
$\beta\le U$, then $y\in B'$.
\end{proposition}

\begin{proof}
Let $x\in B\setminus B'$ and choose $i\in P$ with $x_i\notin B_i'$. Every grid
interval containing $x_i$ is unretained, since retained intervals lie in
$B_i'$. Choose one, $[a,a']$; then $\min\{m_i(a),m_i(a')\}>U$, and
Lemma~\ref{core:lem:interp}(b) gives $F(x)>U$. For the last claim,
$m_i(y_i)\le Q(y)=\beta\le U$, and $y_i$ is an endpoint of at least one grid
interval of coordinate $i$ (each $B_i$ with $i\in P$ has at least two points),
which is therefore retained.
\end{proof}

\begin{definition}[Path certificate]\label{core:def:cert}
A \emph{path certificate} consists of grids $G^{(0)},\dots,G^{(k)}$, a number
$\beta$, and a point $\hat x$. Let $B^{(j)}$ be the subbox spanned by
$G^{(j)}$ (its endpoints are the extreme nodes). It is \emph{valid} if
$\hat x\in X$; $B^{(0)}=X$; $B^{(j+1)}\subseteq B^{(j)}$; and, with $Q^{(j)}$
and $m^{(j)}$ computed from $G^{(j)}$ and the curvature bounds $L_i$:
\begin{enumerate}
\item[(C1)] for $j<k$, every grid interval $[a,a']$ of $G_i^{(j)}$ not
contained in $B_i^{(j+1)}$ has $\min\{m_i^{(j)}(a),m_i^{(j)}(a')\}\ge\beta$;
\item[(C2)] $\min_{G^{(k)}}Q^{(k)}\ge\beta$.
\end{enumerate}
\end{definition}

\begin{theorem}[Path certificates]\label{core:thm:cert}
If a path certificate is valid, then $\beta\le\OPT\le F(\hat x)$. It is
checked by $k+1$ applications of Proposition~\ref{core:prop:dp} and
$O(\sum_{j,i}|G_i^{(j)}|)$ further comparisons, in exact arithmetic, given
the curvature bounds (for a quadratic, the checker computes
$L_i=\max\{0,H_{ii}\}$ itself). No growth or uniqueness assumption is
involved.
\end{theorem}

\begin{proof}
Let $x\in X$ and let $j$ be the largest index with $x\in B^{(j)}$. If $j=k$,
Lemma~\ref{core:lem:interp}(a) and (C2) give $F(x)\ge\beta$. Otherwise some
$x_i\notin B_i^{(j+1)}$. Since $x_i\in B_i^{(j)}$, there are consecutive nodes
$a\le x_i\le a'$ of $G_i^{(j)}$; the interval $[a,a']$ contains $x_i$ and is
therefore not contained in $B_i^{(j+1)}$. By (C1) and
Lemma~\ref{core:lem:interp}(b), $F(x)\ge\beta$.
\end{proof}

\begin{remark}[Minimal content]\label{core:rem:mincert}
Consider any sequence of stages in which $B^{(j+1)}$ is obtained from $G^{(j)}$
by Definition~\ref{core:def:filter} with thresholds $U_j$ that are values of
feasible points, and let $\beta=\beta^{(k)}=\min Q^{(k)}$. Then
$(G^{(0)},\dots,G^{(k)},\beta,\hat x)$ is a valid path certificate for any
feasible $\hat x$. Indeed, every optimizer lies in every $B^{(j)}$
(Proposition~\ref{core:prop:filter}), so $\beta\le\OPT\le U_j$; an interval
not contained in $B_i^{(j+1)}$ is unretained, so its endpoint bound exceeds
$U_j\ge\beta$. Thus a certificate needs only the stage grids (or the boxes
together with the centers, meshes and grading that generate them), one bound
and one feasible point. Messages, min-marginal tables, intermediate incumbents
and unsuccessful trials need not be stored, and stages with
$B^{(j+1)}=B^{(j)}$ may be deleted. The stages at which points are removed
cannot be dropped in general: the last grid bounds $F$ only on $B^{(k)}$.
\end{remark}

\subsection{Graded grids}\label{core:sec:graded}

\begin{definition}[Graded grid]\label{core:def:graded}
Let $B_i=[l_i',u_i']$ (or its integer points), a center $c_i\in B_i$ (an
integer if $i\in I_Z$), a mesh $h>0$, and a grading $0<\theta\le1/4$ be given.
Put $\sigma(t)=h+\theta t$ if $i\in I_C$ and
$\sigma(t)=\max\{1,\lfloor h+\theta t\rfloor\}$ if $i\in I_Z$. Starting from
$t_0=0$, set $t_{k+1}=\min\{t_k+\sigma(t_k),\,u_i'-c_i\}$ until
$t_{k+1}=u_i'-c_i$, and add the nodes $c_i+t_k$; proceed in the same way
towards $l_i'$. The graded grid $G_i$ consists of these nodes.
\end{definition}

\begin{lemma}[Mesh and node count]\label{core:lem:graded}
(a) $w_i(v)\le h+\theta|v-c_i|$ for every $v\in G_i$.
(b) If the distance from $c_i$ to an endpoint is $R>0$, the number of nodes
strictly on that side is at most
$\bigl\lceil\ln(1+\theta R/H)/\ln(1+\theta/3)\bigr\rceil
\le\bigl\lceil(4/\theta)\ln(1+\theta R/H)\bigr\rceil$, where $H=h$ for
$i\in I_C$ and $H=\max\{h,1\}$ for $i\in I_Z$.
(c) If $h\ge u_i'-l_i'$ then $|G_i|\le3$.
\end{lemma}

\begin{proof}
(a) Consider the side towards $u_i'$; the other is symmetric. The interval
between the nodes at distances $t_k<t_{k+1}$ has length at most
$\sigma(t_k)$. For $i\in I_C$ this is $h+\theta t_k$. For $i\in I_Z$ and a
length of at least two, $\sigma(t_k)\ge2$, so
$\sigma(t_k)=\lfloor h+\theta t_k\rfloor\le h+\theta t_k$; intervals of length
one are ignored in $w_i$. The node at distance $t_k$ has adjacent intervals
starting at distances $t_{k-1}$ and $t_k$, both at most $t_k=|v-c_i|$; the
center's adjacent intervals have length at most $h$.

(b) Let $t_1<\dots<t_m=R$ be the nodes on that side. Every step except the
last is unclipped. An unclipped step from $t$ has length
$\sigma(t)\ge(H+\theta t)/3$: for $i\in I_C$ trivially; for $i\in I_Z$ with
$h\ge1$, $\lfloor h+\theta t\rfloor\ge(h+\theta t)/2$ because
$h+\theta t\ge1$; for $i\in I_Z$ with $h<1$, either $\theta t<2$ and
$\sigma\ge1\ge(1+\theta t)/3$, or $\theta t\ge2$ and
$\lfloor h+\theta t\rfloor\ge\theta t/2\ge(1+\theta t)/3$. Hence
$t_{k+1}+H/\theta\ge(1+\theta/3)(t_k+H/\theta)$ for $k\le m-2$, and
$(1+\theta/3)^{m-1}\le(t_{m-1}+H/\theta)\theta/H<1+\theta R/H$. So $m-1$ is
strictly less than $\ln(1+\theta R/H)/\ln(1+\theta/3)$, which gives the first
bound. The second follows from
$\ln(1+x)\ge x-x^2/2\ge\tfrac{23}{24}x$ for $0\le x=\theta/3\le1/12$.

(c) Each side has distance $R\le h$, and $\sigma(0)\ge R$ (for $i\in I_Z$,
$R$ is an integer and $\lfloor h\rfloor\ge R$). The first step reaches the
endpoint.
\end{proof}

\subsection{One trial and its invariants}\label{core:sec:trial}

\begin{definition}[Dyadic meshes and a trial]\label{core:def:trial}
For $i\in P$ let $e_i\in\Z$ satisfy $4^{e_i-1}<L_i\le4^{e_i}$ and put
$r_i=2^{-e_i}$, so that $\frac14<L_ir_i^2\le1$. Let $E$ be the least integer
with $2^Er_i\ge s_i$ for all $i\in P$, and put $\eta_j=2^{E-j}$ and
$h_{ij}=\eta_jr_i$. A \emph{trial} with grading $\theta=2^{-\mu}$, cap $K$,
accuracy $\varepsilon>0$ and stage limit $J$ is the following procedure.
\begin{enumerate}
\item $B^{(0)}=X$, $c^{(0)}=l$ (the lower endpoint vector), and $U$ is the
smallest objective value of a feasible point known so far (at least $F(l)$
is known).
\item For $j=0,\dots,J$: for $i\in P$ let $G_i^{(j)}$ be the graded grid of
$B_i^{(j)}$ with center $c_i^{(j)}$, mesh $h_{ij}$ and grading $\theta$; for
$i\notin P$ let $G_i^{(j)}=\{l_i,u_i\}$. If some $|G_i^{(j)}|$ would exceed
$K$, stop generating and \emph{abort} the trial before forming any table.
Otherwise compute $\beta_j$, a minimizer $y^{(j)}$ and all $m_i$
(Proposition~\ref{core:prop:dp}); set $U\leftarrow\min\{U,F(y^{(j)})\}$. If
$U-\beta_j\le\varepsilon$, stop with \emph{success}. Otherwise filter with
threshold $U$ to obtain $B^{(j+1)}$ and set $c^{(j+1)}=y^{(j)}$.
\end{enumerate}
\end{definition}

The centers are feasible: $c^{(0)}=l$, and $y^{(j)}\in B^{(j+1)}$ by
Proposition~\ref{core:prop:filter}, because $\beta_j\le\min_{B^{(j)}}F\le U$
(the incumbent lies in $B^{(j)}$ by the same proposition).

\begin{lemma}[Stage invariants]\label{core:lem:inv}
Assume weighted growth with constant $\gamma$ and minimizer $x^*$, let
$\bar\kappa\ge\max\{1,1/\gamma\}$ and $8\bar\kappa\theta^2\le1$, and write
$a_j=n_P\eta_j^2$. At every stage $j$ reached by the trial:
\begin{enumerate}
\item[(i)] $x^*\in B^{(j)}$ and $\beta_j\le\OPT$;
\item[(ii)] $\norm{c^{(j)}-x^*}_{\mathbf L}^2\le4\bar\kappa a_j$;
\item[(iii)] $\norm{y^{(j)}-x^*}_{\mathbf L}^2\le\bar\kappa a_j$;
\item[(iv)] $U-\beta_j\le F(y^{(j)})-\beta_j=D(y^{(j)})\le\frac9{16}a_j$;
\item[(v)] every $z\in G^{(j)}$ with $Q^{(j)}(z)\le U$ (after the update at
stage $j$) satisfies $\norm{z-x^*}_{\mathbf L}^2\le\frac{17}{15}\bar\kappa a_j$.
\end{enumerate}
\end{lemma}

\begin{proof}
(i) follows from Proposition~\ref{core:prop:filter}, since $U\ge\OPT$ at every
stage, and from Lemma~\ref{core:lem:interp}(a).

We first show that every $z\in G^{(j)}$ satisfies, with $c=c^{(j)}$,
\begin{equation}\label{core:eq:energy}
D(z)\le\frac{a_j}4+\frac{\theta^2}2\Bigl(\norm{z-x^*}_{\mathbf L}^2
+\norm{c-x^*}_{\mathbf L}^2\Bigr).
\end{equation}
For $i\in P$, Lemma~\ref{core:lem:graded}(a) and $L_ih_{ij}^2\le\eta_j^2$ give
$d_i(z_i)\le\frac{L_i}8(h_{ij}+\theta|z_i-c_i|)^2
\le\frac{\eta_j^2}4+\frac{\theta^2}4L_i(z_i-c_i)^2$, and
$(z_i-c_i)^2\le2(z_i-x_i^*)^2+2(c_i-x_i^*)^2$. For $i\notin P$, $d_i=0$.

(ii) For $j=0$, $\norm{l-x^*}_{\mathbf L}^2\le\sum_{i\in P}L_is_i^2\le
\sum_{i\in P}L_ir_i^2\eta_0^2\le a_0$. For $j\ge1$ it follows from (iii) at
stage $j-1$, because $a_{j-1}=4a_j$.

(iii) Write $y=y^{(j)}$, $Y=\norm{y-x^*}_{\mathbf L}^2$ and $a=a_j$. Since
$\theta^2/2\le1/(16\bar\kappa)\le\gamma/16$, growth, (i),
Lemma~\ref{core:lem:interp}(c), \eqref{core:eq:energy} and (ii) give
\[
\gamma Y\le F(y)-\OPT\le F(y)-\beta_j=D(y)\le\frac a4+\frac\gamma{16}(Y+4\bar\kappa a),
\]
so $Y\le\frac4{15}(\gamma^{-1}+\bar\kappa)a\le\frac8{15}\bar\kappa a$.

(iv) By \eqref{core:eq:energy}, (ii), (iii) and $\theta^2\bar\kappa\le1/8$,
$D(y)\le\frac a4+\frac{\theta^2}2\cdot5\bar\kappa a\le\frac9{16}a$; also
$U\le F(y)$.

(v) Let $Z=\norm{z-x^*}_{\mathbf L}^2$. Using $Q(z)\le U\le F(y)$,
$\beta_j\le\OPT$, (iv) and \eqref{core:eq:energy},
\[
\gamma Z\le F(z)-\OPT=Q(z)+D(z)-\OPT\le D(y)+D(z)
\le\frac{13}{16}a+\frac\gamma{16}(Z+4\bar\kappa a),
\]
so $Z\le(\frac{13}{15}\gamma^{-1}+\frac4{15}\bar\kappa)a\le\frac{17}{15}\bar\kappa a$.
\end{proof}

\begin{lemma}[Localization and uniform state bound]\label{core:lem:states}
Under the hypotheses of Lemma~\ref{core:lem:inv}, for every stage $j$ that is
followed by filtering and every $i\in P$,
\[
B_i^{(j+1)}\subseteq\bigl[y_i^{(j)}-R_{ij},\,y_i^{(j)}+R_{ij}\bigr],\qquad
R_{ij}=8\sqrt{\bar\kappa n_P}\,h_{ij}+[i\in I_Z].
\]
Consequently every grid of the trial has at most
\begin{equation}\label{core:eq:cap}
K(\theta,n_P)=10\,\theta^{-1}\lceil\log_2(n_P+2)\rceil
\end{equation}
nodes per coordinate, a number independent of the stage index and of
$\varepsilon$.
\end{lemma}

\begin{proof}
Write $y=y^{(j)}$, $c=c^{(j)}$, $h=h_{ij}$, $\rho=\sqrt{\bar\kappa n_P}$. Let
$[v,v']$ be a retained interval of coordinate $i$ whose endpoint $v$ satisfies
$m_i(v)\le U$, and let $z\in G^{(j)}$ attain $m_i(v)$, so $z_i=v$. Because
$\sqrt{L_i}>1/(2r_i)$, a bound $\sqrt{L_i}|\delta|\le\lambda\rho\eta_j$
implies $|\delta|<2\lambda\rho h$. Lemma~\ref{core:lem:inv}(ii), (iii), (v)
therefore give
\[
|v-y_i|<2\bigl(1+\sqrt{17/15}\bigr)\rho h,\qquad
|v-c_i|<2\bigl(2+\sqrt{17/15}\bigr)\rho h.
\]
If $[v,v']$ is an integer interval of length one, every point of it is within
$|v-y_i|+1\le1+8\rho h$ of $y_i$. Otherwise its length is at most
$w_i(v)\le h+\theta|v-c_i|$ (Lemma~\ref{core:lem:graded}(a)), and
$\theta\sqrt{\bar\kappa}\le1/\sqrt8$ gives
$\theta|v-c_i|<(2+\sqrt{17/15})\sqrt{n_P/2}\,h$. Since $h\le\rho h$ and
$\sqrt{n_P}\le\rho$, every point of the interval is within
\[
\Bigl(2+2\sqrt{17/15}+1+\bigl(2+\sqrt{17/15}\bigr)/\sqrt2\Bigr)\rho h
<7.3\,\rho h
\]
of $y_i$. The hull of the retained intervals lies in the same interval.

At stage $j+1$ the grid of coordinate $i\in P$ is centered at $y_i$ with mesh
$h/2$, and each side has length at most $R_{ij}$. With $H$ as in
Lemma~\ref{core:lem:graded}(b) for the mesh $h/2$, we have $H\ge h/2$, and
$H\ge1$ for $i\in I_Z$, so
$\theta R_{ij}/H\le\theta+16\theta\rho\le\frac14+4\sqrt{2n_P}$. By
Lemma~\ref{core:lem:graded}(b),
\[
|G_i^{(j+1)}|\le1+2\Bigl\lceil\frac4\theta\ln\bigl(\tfrac54+4\sqrt{2n_P}\bigr)\Bigr\rceil
\le\frac1\theta\Bigl(\frac34+8\ln\bigl(\tfrac54+4\sqrt{2n_P}\bigr)\Bigr),
\]
using $3\le3/(4\theta)$. Call the bracket $\varphi(n_P)$. Then
$\varphi(1)<16.3\le10\lceil\log_23\rceil$, $\varphi(2)<18.6\le10\log_24$, and
for $m\ge2$, $\varphi'(m)\le4/m$ while
$\frac{d}{dm}10\log_2(m+2)=\frac{10}{(m+2)\ln2}\ge\frac{7.2}m$; hence
$\varphi(m)\le10\log_2(m+2)$ for all $m\ge2$. Stage $0$ has at most three nodes
per coordinate by Lemma~\ref{core:lem:graded}(c), since $h_{i0}=2^Er_i\ge s_i$;
coordinates outside $P$ have at most two.
\end{proof}

With a common bound $L_i=L$ and the undyadic meshes $h_j=s2^{-j}$ the same
proof gives the radius $4.2\sqrt{\kappa n}\,h_j$ and the cap
$8\theta^{-1}\lceil\log_2(n+2)\rceil$.

\begin{remark}\label{core:rem:cluster}
The mechanism of Lemmas~\ref{core:lem:inv} and~\ref{core:lem:states} is the
one behind the analysis of the cluster problem in branch and bound
\cite{DuKearfott1994,WechsungSchaberBarton2014}: a lower bound whose error is
quadratic in the cell width, combined with quadratic growth, leaves a number of
cells near the minimizer that does not depend on the accuracy, provided the
error prefactor is small relative to the growth. Here the prefactor condition
is $8\bar\kappa\theta^2\le1$. The differences are that growth is global,
the cells form product grids handled exactly by tree dynamic programming, and
the resulting bound is turned into a bit-complexity statement.
\end{remark}

\subsection{The capped schedule without a growth constant}\label{core:sec:schedule}

\begin{definition}[Schedule]\label{core:def:schedule}
If $P=\emptyset$, run one stage with $G_i=\{l_i,u_i\}$ for all $i$. Otherwise
let $J$ be the least integer $J\ge0$ with
$\frac9{16}n_P\eta_0^24^{-J}\le\varepsilon$, and for $\mu=2,3,\dots$ run a
trial with $\theta=2^{-\mu}$, cap $K_\mu=10\cdot2^\mu\lceil\log_2(n_P+2)\rceil$
and stage limit $J$, until a trial succeeds. Return the final incumbent
$\hat x$, the bound $\beta$ of the successful stage, and the grids of the
successful trial.
\end{definition}

\begin{theorem}[Validity, termination, and work]\label{core:thm:schedule}
For every instance and every $\varepsilon>0$:
\begin{enumerate}
\item[(a)] The schedule halts and returns a feasible $\hat x$ and a valid path
certificate with $\beta\le\OPT\le F(\hat x)\le\beta+\varepsilon$.
\item[(b)] If weighted growth holds, then with
$\mu^*=\max\{2,\lceil\log_4(8\bar\kappa)\rceil\}$ the schedule succeeds at
some trial $\mu\le\mu^*$, $2^{\mu^*}\le6\sqrt{\bar\kappa}$, and the total
number of table operations is at most
\[
c\,(J+1)\,p\,(|\mathcal A|+n+N)\,K_{\mu^*}^p,
\]
for an absolute constant $c$, plus the factor evaluations at grid points.
\end{enumerate}
\end{theorem}

\begin{proof}
If $P=\emptyset$, then $D\equiv0$ and $F(y)=\beta\le\OPT$, so the single stage
is exact. Assume $P\ne\emptyset$.

(a) Validity of every stage bound and of the certificate is
Lemma~\ref{core:lem:interp}, Proposition~\ref{core:prop:filter} and
Remark~\ref{core:rem:mincert}. For termination, choose $\mu\ge2$ with
$\theta=2^{-\mu}\le\min_{i\in P}h_{iJ}/s_i$. In this trial, at every stage
$j\le J$ and for $i\in P$ we have $\theta s_i\le h_{ij}$. Continuous steps are
at least $h_{ij}$, so $|G_i^{(j)}|\le3+s_i/h_{ij}\le3+1/\theta$. Integer steps
are at least $\lfloor h_{ij}\rfloor\ge h_{ij}/2$ if $h_{ij}\ge1$, giving
$|G_i^{(j)}|\le3+2/\theta$; if $h_{ij}<1$ then
$|G_i^{(j)}|\le s_i+1<1/\theta+1$. All these are below $K_\mu\ge20/\theta$, so
the trial is not aborted. At stage $J$, $w_i\le h_{iJ}+\theta s_i\le2h_{iJ}$,
so $U-\beta_J\le D(y^{(J)})\le\sum_{i\in P}L_ih_{iJ}^2/2\le n_P\eta_J^2/2
\le\frac89\varepsilon$, and the trial succeeds unless an earlier one did.

(b) The choice of $\mu^*$ gives $8\bar\kappa\theta^2\le1$. By
Lemma~\ref{core:lem:states} that trial is never aborted, and by
Lemma~\ref{core:lem:inv}(iv) it reaches $U-\beta_J\le\frac9{16}a_J\le\varepsilon$.
If $\mu^*=2$ then $2^{\mu^*}=4\le6\sqrt{\bar\kappa}$; otherwise
$4^{\mu^*-1}<8\bar\kappa$ and $2^{\mu^*}<2\sqrt8\sqrt{\bar\kappa}$. Each trial
$\mu\le\mu^*$ runs at most $J+1$ stages with grids of at most $K_\mu$ nodes
(the aborted stage only generates nodes), so by
Proposition~\ref{core:prop:dp} it costs $O((J+1)p(|\mathcal A|+n+N)K_\mu^p)$
table operations; $\sum_{\mu\le\mu^*}K_\mu^p\le2K_{\mu^*}^p$.
\end{proof}

The caps are essential to (b): without them an unsuitable early trial could
retain a large part of the box at every stage and pay a power of the accuracy
before an admissible trial is reached.

\subsection{Bit complexity}\label{core:sec:bits}

\begin{theorem}[Fixed-parameter bit complexity]\label{core:thm:bits}
Let the factors be explicit rational quadratics given as monomial lists, let
$I$ be the binary input length (including the decomposition), let
$L_i=\max\{0,\partial_{ii}^2F\}$, and let $\varepsilon=2^{-q}$. The schedule
always halts with a valid certificate. If weighted growth holds, it uses at
most
\[
f(p,\bar\kappa)\,(I+q+1)^5,\qquad
f(p,\kappa)=c_0\,(c_1p\sqrt\kappa)^p\,\kappa\,(1+\log_2\kappa)^2,
\]
bit operations, for absolute constants $c_0,c_1$; the certificate has size,
and can be checked in time, within the same bound. In particular the bound
holds with $\kappa\ge\bar\kappa$ in place of $\bar\kappa$.
\end{theorem}

\begin{proof}
Let $\Gamma_X$ be the product of the denominators of the endpoints and
$\Gamma_F$ the product of the denominators of the coefficients, after
substituting fixed coordinates; both have $O(I)$ bits, and
$|e_i|,|E|=O(I)$. Fix a trial $\mu$ and let $\alpha=J+\max_i|e_i|+|E|+\mu K_\mu$.

\emph{Denominators.} We claim that every node of every grid of the trial has
a denominator dividing $\Gamma_X2^\alpha$. The initial center and the original
endpoints have denominators dividing $\Gamma_X$. A new node is a center
coordinate plus or minus
$h_{ij}\frac{(1+\theta)^k-1}\theta
=2^{E-j-e_i}\frac{(2^\mu+1)^k-2^{\mu k}}{2^{\mu(k-1)}}$ with
$k\le K_\mu+1$ (one extra node detects an abort), or an integer step for
$i\in I_Z$, or a clipped endpoint of the current box. Centers and box
endpoints are nodes of earlier stages of the trial or original data. The
claim follows by induction, and denominators do not grow with the number of
stages.

\emph{Bit lengths.} Nodes lie in $\bar X$, so their numerators have $O(I+\alpha)$
bits. Every factor value, correction $L_iw^2/8$, table entry and message has
a denominator dividing $8\Gamma_F\Gamma_X^24^\alpha$ and is a sum of at most
$|\mathcal A|+n$ values of the form $\pm$coefficient $\times$ product of at
most two nodes (a message is the value of such a sum at a minimizing
assignment). Writing all of them over this common denominator, every number
is an integer of $b=O(I+\alpha)$ bits, the dynamic program uses only integer
additions and comparisons, and a factor evaluation costs $O(b^2)$ bit
operations per monomial.

\emph{Counting.} We have $J+1\le c(I+q+1)$, and
$\ell:=\lceil\log_2(n_P+2)\rceil\le I+1$, so $\mu K_\mu\le20\mu2^\mu(I+1)$ and
$b\le c(1+\mu2^\mu)(I+q+1)$. Since $\sup_{t\ge0}t^pe^{-t}=(p/e)^p$ and
$\ell\le\ln(2(n+2))/\ln2$,
\[
\ell^p\le\Bigl(\frac{p}{e\ln2}\Bigr)^p2(n+2)\le2p^p(n+2)\le6p^pI,
\]
so $K_\mu^p\le6I(10\cdot2^\mu p)^p$. By Theorem~\ref{core:thm:schedule} and
Proposition~\ref{core:prop:dp}, trial $\mu$ costs at most
$c(J+1)\,pI\,K_\mu^p\,b^2\le c'\,p(10\cdot2^\mu p)^p(1+\mu2^\mu)^2(I+q+1)^5$ bit
operations, including grid generation and factor evaluation; preprocessing
(computing $L_i$, $e_i$, $E$, $J$ by rational comparisons) is polynomial and of
lower order. Summing over $\mu\le\mu^*$ at most doubles the last term, and
$2^{\mu^*}\le6\sqrt{\bar\kappa}$, $\mu^*\le3(1+\log_2\bar\kappa)$ give the
stated $f$. The certificate stores the grids of one trial and is checked by
the same dynamic programs. Without growth, the schedule still halts by
Theorem~\ref{core:thm:schedule}(a), and every number it forms is a rational
whose size is bounded as above for its trial.
\end{proof}

\begin{remark}[In what sense the result is fixed-parameter tractable]
\label{core:rem:fpt}
(i) The parameter is $(p,\bar\kappa)$. The running time is at most
$f(p,\bar\kappa)(I+q+1)^5$ for every instance with weighted growth, where
$\bar\kappa$ is a property of the instance that the algorithm neither
receives nor computes; $f$ is computable and the exponent $5$ does not depend
on $p$ or $\bar\kappa$. The accuracy enters through $q=\log_2(1/\varepsilon)$
only polynomially. No bound on variable occurrences, on Hessian norms (other
than through $\bar\kappa$), or on the number of local minima is used.
(ii) Validity of the output never depends on growth, uniqueness, or the trial
at which the schedule stops; these affect only the running-time bound.
Instances without growth (for example with several minimizers that differ in
$P$) are solved to accuracy $\varepsilon$ in finite time, without the bound.
(iii) The bag size $p$ refers to the supplied decomposition. Since every
factor scope is a clique of the primal graph, a tree decomposition of width
at most $2\,\mathrm{tw}+1$, computable in time $2^{O(\mathrm{tw})}n$
\cite{Korhonen2021}, can be used instead; the bound then holds with
$p=2\,\mathrm{tw}+2$.
(iv) The exponent $5$ reflects schoolbook multiplication and is not
optimized.
\end{remark}

\subsection{Both parameters are necessary}\label{core:sec:lower}

The following statements concern any algorithm, not only the one above. ETH
denotes the exponential-time hypothesis \cite{ImpagliazzoPaturiZane2001};
rETH denotes its randomized version (no randomized algorithm with bounded
error solves 3-SAT in time $2^{o(n)}$).

\begin{proposition}[Exponential dependence on the width at bounded conditioning]
\label{core:prop:lbwidth}
For a graph $(V,E)$ with $V=\{1,\dots,n\}$ let
\[
\Psi(x)=2^n\Bigl(-\sum_ix_i+2\sum_{ij\in E}x_ix_j\Bigr)+\sum_i2^{i-1}x_i
+\frac1{2n}\sum_i(x_i^2-x_i),\qquad x\in[0,1]^n .
\]
Then $\Psi$ has a unique minimizer $x^*$, which is binary,
$\lfloor\Psi^*/2^n\rfloor=-\alpha(V,E)$ (the independence number), and
$\Psi(x)-\Psi^*\ge\frac1{2n}\norm{x-x^*}^2$ on $[0,1]^n$. With
$L_i=L=1/n$ this gives $\kappa\le2$ (and $\bar\kappa\le2$). The same holds on
$\{0,1\}^n$. Every tree decomposition of the graph is a decomposition of
$\Psi$. Consequently, assuming ETH, no algorithm computes $\OPT$ within $1/2$
for all box quadratics with $\kappa\le2$ (continuous, or all-binary) in time
$2^{o(p)}I^{O(1)}$.
\end{proposition}

\begin{proof}
On binary points $\Psi=2^n\Phi+\tau$ with $\Phi=-\sum x_i+2\sum_Ex_ix_j$
integral and $\tau=\sum2^{i-1}x_i\in\{0,\dots,2^n-1\}$ injective. A binary
minimizer therefore minimizes $\Phi$ and is unique. If the support $S$ of a
binary $x$ spans an edge, deleting an endpoint $u\in S$ of such an edge changes
$\Phi$ by $1-2\deg_S(u)\le-1$; hence $\min\Phi=-\alpha$ and
$\lfloor\Psi^*/2^n\rfloor=-\alpha$.

For $x\in[0,1]^n$ write $\Psi=M+\frac1{2n}\sum(x_i^2-x_i)$, where $M$ is
multilinear. Let $V$ have independent coordinates with $\Pr[V_i=1]=x_i$. Then
$M(x)=\E M(V)=\E\Psi(V)$, and since $\Psi(v)\ge\Psi^*+1$ for binary
$v\ne x^*$, $M(x)-\Psi^*\ge\Pr[V\ne x^*]$. With
$\delta_i=|x_i-x_i^*|\in[0,1]$,
$\Pr[V=x^*]=\prod_i(1-\delta_i)\le1-\max_i\delta_i$, so
$M(x)-\Psi^*\ge\max_i\delta_i\ge\frac1n\sum_i\delta_i$. Also
$x_i(1-x_i)\le\delta_i$ because $x_i^*\in\{0,1\}$. Hence
$\Psi(x)-\Psi^*\ge(\frac1n-\frac1{2n})\sum_i\delta_i\ge\frac1{2n}\norm{x-x^*}^2$.
The only squared terms are $\frac1{2n}x_i^2$, so $L=1/n$ and $L/g\le2$.

The instance has $O(n^2)$ coefficients of $O(n)$ bits, so $I=n^{O(1)}$, and
the single bag $V$ gives $p=n$. An algorithm as excluded would compute
$\alpha$ in time $2^{o(n)}$, contradicting ETH
\cite{ImpagliazzoPaturiZane2001}.
\end{proof}

Under the strong exponential-time hypothesis the same reduction, applied with
a supplied tree decomposition of the graph, excludes $(2-\epsilon)^pI^{O(1)}$
by the lower bound of \cite{LokshtanovMarxSaurabh2018} for independent set.

\begin{proposition}[Polynomial dependence on conditioning at width three]
\label{core:prop:lbkappa}
Given positive integers $a_1,\dots,a_m$ and $t\ge0$, let $A=\sum_ka_k$,
$M=(A+t)^2+1$, and on $\sigma\in\{0,1\}^m$, $s\in(\{0,\dots,A\}\cap\Z)^m$, with
$s_0=0$,
\[
\Psi=2^{m+1}\Bigl(M\sum_{k=1}^m(s_k-s_{k-1}-a_k\sigma_k)^2+(s_m-t)^2\Bigr)
+\sum_k2^{k-1}\sigma_k .
\]
The bags $\{s_{k-1},\sigma_k,s_k\}$ form a path decomposition with $p=3$.
$\Psi$ has a unique minimizer, $\Psi(z)\ge\Psi^*+1$ for every other $z$,
$\Psi^*<2^{m+1}$ if and only if some subset of the $a_k$ sums to $t$, and
\[
\log_2\kappa\le m+O(\log(m+A+t)).
\]
Consequently: unless $\mathrm{P}=\mathrm{NP}$, no algorithm computes $\OPT$
within $1/2$ for all integer box quadratics with $p\le3$ and finite $\kappa$
in time polynomial in $I$ and $\log\kappa$; and assuming ETH there is a
constant $c>0$ such that none does so in time $\kappa^cI^{O(1)}$.
\end{proposition}

\begin{proof}
Let $\Phi$ be the bracket. For fixed $\sigma$, the partial sums
$s^\sigma_k=\sum_{l\le k}a_l\sigma_l\in[0,A]$ give
$\Phi(\sigma,s^\sigma)=(s^\sigma_m-t)^2\le(A+t)^2<M$, while every other $s$
violates a recursion and has $\Phi\ge M$; so $s^\sigma$ is the unique minimizer
of $\Phi(\sigma,\cdot)$. Since $\Phi$ is an integer and the tie-break term lies
in $[0,2^m-1]$ and is injective in $\sigma$, $\Psi$ has a unique minimizer and
integral values, which gives the gap. $\Phi=0$ exactly at $(\sigma,s^\sigma)$
with $s^\sigma_m=t$, and $\Phi\ge1$ implies $\Psi\ge2^{m+1}$.

Every coordinate ranges over an interval of length at most $A$, so
$\norm{z-z^*}^2\le m(1+A^2)$ and $g\ge1/(m(1+A^2))$. The diagonal Hessian
entries are $2^{m+2}Ma_k^2$, $2^{m+3}M$ and $2^{m+2}(M+1)$, so
$L\le2^{m+3}M(\max_ka_k^2+1)$ and the bound on $\kappa\le L/g$ follows. The
input length is polynomial in $m$ and $\log(A+t)$.

An algorithm polynomial in $I$ and $\log\kappa$ would therefore decide Subset
Sum in polynomial time. Under ETH, Subset Sum with $m$ numbers of $O(m)$ bits
has no $2^{o(m)}$ algorithm (from the sparsification lemma
\cite{ImpagliazzoPaturiZane2001} and the textbook reduction from 3-SAT, which
produces $O(n)$ numbers of $O(n)$ bits after sparsification). For such inputs
$\log_2\kappa\le C_1m$ and $I=O(m^2)$, so a running time $\kappa^cI^{O(1)}$
with $c$ below the ETH constant divided by $C_1$ is excluded.
\end{proof}

\begin{lemma}[Isolation \cite{MulmuleyVaziraniVazirani1987}]\label{core:lem:isolation}
Let $\mathcal F$ be a nonempty family of subsets of a finite set $W$, and let
$w:W\to\{1,\dots,R\}$ be uniformly random. Then the minimum of
$w(S)=\sum_{e\in S}w(e)$ over $S\in\mathcal F$ is attained by a unique set with
probability at least $1-|W|/R$.
\end{lemma}

\begin{proof}
For $e\in W$ let $\alpha_e=\min\{w(S):S\in\mathcal F,e\notin S\}
-\min\{w(S\setminus\{e\}):S\in\mathcal F,e\in S\}$ (with
$\min\emptyset=+\infty$); $\alpha_e$ does not depend on $w(e)$. If two
distinct minimizing sets exist, some $e$ lies in one and not in the other,
and then $w(e)=\alpha_e$, an event of probability at most $1/R$.
\end{proof}

\begin{proposition}[The exponent of the conditioning must grow with the width]
\label{core:prop:lbproduct}
Assume rETH. There are no computable $f$ and computable nondecreasing
unbounded $\psi$ such that some algorithm, on every integer box quadratic with
supplied decomposition and a unique minimizer, outputs a feasible point within
$1/2$ of $\OPT$ in time $f(p)\,(\kappa I)^{p/\psi(p)}$. In particular a bound
$f(p)\kappa^{a(p)}I^{C}$ requires $a(p)\ne o(p)$; the bound of
Theorem~\ref{core:thm:bits} has $a(p)=p/2+O(1)$.
\end{proposition}

\begin{proof}
Take an instance of multicolored clique: $k$ color classes identified with
$\{1,\dots,N\}$ and, for $c<d$, edge lists
$R_{cd}=((a_1,b_1),\dots,(a_m,b_m))$ of distinct pairs. Choose independent
uniform weights $w(e)\in\{1,\dots,2|R|\}$ on the edges $e$ of
$R=\bigsqcup R_{cd}$, and let $W_0=1+\binom k2\,2|R|$. Use integer variables
$x_c\in[1,N]$ and, for each pair $c<d$ and $l=1,\dots,m=|R_{cd}|$,
$\zeta_l,S_l\in\{0,1\}$ and $A_l,B_l\in[0,N]$, with $S_0=A_0=B_0=0$. Let
\begin{align*}
\Phi&=\sum_{c<d}\Bigl[\sum_{l}\bigl((S_l-S_{l-1}-\zeta_l)^2
+(A_l-A_{l-1}-a_l\zeta_l)^2+(B_l-B_{l-1}-b_l\zeta_l)^2\bigr)\\
&\qquad\qquad+(S_m-1)^2+(A_m-x_c)^2+(B_m-x_d)^2\Bigr],\qquad
\Psi_w=W_0\Phi+\sum_{c<d}\sum_lw(a_l,b_l)\zeta_l .
\end{align*}
\emph{Decomposition.} Take a central bag $\{x_1,\dots,x_k\}$; for each pair a
bag $\{x_c,x_d,S_m,A_m,B_m\}$ adjacent to it, followed by the path of bags
$\{S_{l-1},A_{l-1},B_{l-1},\zeta_l,S_l,A_l,B_l\}$, $l=m,\dots,1$. Every term
lies in a bag, running intersection holds, and $p=\max\{k,7\}$.

\emph{Encoding.} $\Phi\ge0$ is an integer. $\Phi=0$ forces $S$ to increase
from $0$ to $1$ in steps $\zeta_l\in\{0,1\}$, so exactly one $\zeta_{l^*}=1$,
and then $(x_c,x_d)=(A_m,B_m)=(a_{l^*},b_{l^*})\in R_{cd}$. Thus $\Phi=0$
exactly at encodings of multicolored cliques, and each clique has exactly one
encoding because the pairs in $R_{cd}$ are distinct. At an encoding the
tie-break equals the weight of the clique's edge set. If a clique exists and
the minimum-weight clique is unique, $\Psi_w$ has a unique minimizer, value at
most $W_0-1$, and integral values; every point with $\Phi\ge1$ has
$\Psi_w\ge W_0$.

\emph{Conditioning.} All domains have length at most $N$, and there are
$n_v\le k+2k^2N^2$ variables, so $g\ge1/(n_vN^2)$; the diagonal Hessian entries
are at most $W_0(2+4N^2+2k)$. Hence $\kappa\le W_0(2+4N^2+2k)n_vN^2$, and both
$\kappa$ and $I$ are at most $(kN)^{c_2}$ for an absolute $c_2$.

\emph{Reduction.} Run the assumed algorithm for
$f(p)((kN)^{2c_2})^{p/\psi(p)}$ steps. If it outputs a feasible $\hat z$ with
$\Psi_w(\hat z)<W_0$, then $\Phi(\hat z)=0$ and $\hat z$ encodes a clique; answer
yes. Otherwise answer no. There are no false positives. If a clique exists,
Lemma~\ref{core:lem:isolation} (with $W=R$ and $\mathcal F$ the edge sets of
cliques) makes the minimum-weight clique unique with probability at least
$1/2$, and then the algorithm halts in time with $\Psi_w(\hat z)\le W_0-1/2$.
This is a one-sided-error algorithm for multicolored clique with running time
$f'(k)N^{o(k)}$, which rETH excludes \cite{ChenEtAl2006} (the reductions
behind that bound are deterministic, so they also exclude randomized
algorithms under rETH).
\end{proof}

\begin{remark}
Propositions~\ref{core:prop:lbwidth}--\ref{core:prop:lbproduct} show that
neither parameter can be dropped and that the shape $f(p)\kappa^{\Theta(p)}$
cannot be improved to $f(p)\kappa^{o(p)}$ in general. They leave open the
constant in the exponent, whether the factor $p^p$ (coming from
$\lceil\log_2(n+2)\rceil^p$) is necessary, and whether
Proposition~\ref{core:prop:lbproduct} holds for purely continuous instances
or with a deterministic reduction.
\end{remark}

\subsection{A nonconvex family covered by the theorem}\label{core:sec:example}

\begin{example}\label{core:ex:family}
For a block $b$ with variables $(u_b,v_b,\rho_b)\in[0,1]^3$ let
\[
\phi(u,v,\rho)=u^2+v^2-4uv+\tfrac14(u+v)+(\rho-\tfrac u2)^2 .
\]
For a graph $\Gamma$ on the blocks with maximum degree $\Delta$ let
\[
F=\sum_b\phi(u_b,v_b,\rho_b)+\frac1{16\Delta}\sum_{bb'\in\Gamma}(u_b-u_{b'})^2 .
\]
Then: (a) $x^*=(1,1,\frac12)$ in every block is the unique minimizer and
$F(x)-\OPT\ge\frac12\norm{x-x^*}^2$; (b) $L\le\frac{21}8$, so
$\kappa\le\frac{21}4$; (c) each of the $2^m$ points with
$(u_b,v_b,\rho_b)\in\{(0,0,0),(1,1,\frac12)\}$ for all $m$ blocks is a strict
local minimizer; (d) a tree decomposition of $\Gamma$ of width $w$ yields one
of $F$ with $p=\max\{w+1,3\}$. Each block Hessian is indefinite (along
$(1,1,\frac12)$ the second-order term of $\phi$ equals $-2$), and the
minimizer is not a vertex. For $\Gamma$ a $k\times m'$ grid with $m'\ge k$,
Theorem~\ref{core:thm:bits} solves this family in time
$f(\max\{k+1,3\},\frac{21}4)(I+q+1)^5$, although it has $2^{km'}$ strict
local minima.
\end{example}

\begin{proof}
(a) Let $d=(1-u,1-v)\in[0,1]^2$. Since $\nabla_{u,v}(u^2+v^2-4uv+\frac{u+v}4)
=(-\frac74,-\frac74)$ at $(1,1)$ and the quadratic part is $d_1^2+d_2^2-4d_1d_2
\ge-(d_1^2+d_2^2)$,
\[
u^2+v^2-4uv+\tfrac{u+v}4+\tfrac32\ge\tfrac74(d_1+d_2)-(d_1^2+d_2^2)
\ge\tfrac34(d_1^2+d_2^2),
\]
using $d_i\ge d_i^2$. With $(\rho-\frac12)^2\le2(\rho-\frac u2)^2+\frac12d_1^2$,
\[
\phi-\phi^*\ge\tfrac34(d_1^2+d_2^2)+(\rho-\tfrac u2)^2
\ge\tfrac12\bigl(d_1^2+d_2^2+(\rho-\tfrac12)^2\bigr).
\]
The coupling is nonnegative and vanishes at $x^*$. Summing over blocks gives
(a). (b) The diagonal entries are $\frac52+\frac{2\deg(b)}{16\Delta}$ for
$u_b$ and $2$ for $v_b,\rho_b$.

(c) At such a point $\bar x$, the coupling gradient in $u_b$ has absolute
value at most $\frac{2\Delta}{16\Delta}=\frac18$. Hence $\partial_{u_b}F\ge\frac18$
and $\partial_{v_b}F=\frac14$ at $u_b=v_b=0$, $\partial_{u_b}F\le-\frac74+\frac18$
and $\partial_{v_b}F=-\frac74$ at $u_b=v_b=1$, and $\partial_{\rho_b}F=0$. For a
feasible $x$ let $d=x-\bar x$ and $\delta=\sum_b(|d_{u_b}|+|d_{v_b}|)$. The
first-order term is at least $\delta/8$. The second-order term equals
$\sum_b[(d_{\rho_b}-\frac12d_{u_b})^2+d_{u_b}^2+d_{v_b}^2-4d_{u_b}d_{v_b}]$ plus a
nonnegative coupling, which is at least
$\sum_b(d_{\rho_b}-\frac12d_{u_b})^2-\delta^2$. Thus
$F(x)-F(\bar x)\ge\delta/8-\delta^2+\sum_b(d_{\rho_b}-\frac12d_{u_b})^2>0$ when
$0<\delta<\frac18$, and $=\sum_bd_{\rho_b}^2>0$ when $\delta=0\ne d$.
(d) Attach to a bag containing $u_b$ a new bag $\{u_b,v_b,\rho_b\}$.
\end{proof}
